\section{Normalization of continuous-arc drawings}
\label{sec:normalization}

This appendix supplies the drawing reduction used in
Section~\ref{sec:intersections}.  The argument applies to any finite
simple graph and counts every crossing point separately. The removal
of adjacent-edge crossings is a standard ordinary-crossing reduction
\cite{FulekEtAl2012}; we include the continuous-arc details.

\begin{lemma}[Drawing normalization]
\label{lem:normalization}
An admissible continuous-arc drawing of a finite simple graph can be
replaced by a polygonal drawing with no more crossing points, with
edges straight near every vertex and crossing, and with no crossings
between edges sharing an endpoint.  In particular, a drawing attaining
the ordinary crossing number can be chosen to have these properties.
\end{lemma}

\begin{proof}
\emph{Isolating the marked passages.}
Call the vertices and crossing points of the given drawing its
\emph{marks}.  There are finitely many.  Parameterize each edge
injectively by $\gamma_e:[0,1]\to\R^2$.  On each edge, choose pairwise
disjoint parameter neighborhoods of its marked parameters, with
disjoint closures, using one-sided neighborhoods at $0$ and $1$.
For a mark $w$ on that edge, the image of the compact parameter set
outside its chosen neighborhood avoids $w$, and hence has positive
distance from $w$.  If $w$ is not on the edge, the entire compact edge
image has positive distance from $w$.  Since all these conditions are
finite, we may choose pairwise disjoint closed round disks $D_w$,
centered at the marks, so that every visit of an edge to $D_w$ is in
its prescribed parameter neighborhood, and nonincident edges avoid
$D_w$ altogether.

For an interior mark on $e$, remove the parameter interval from the
first to the last visit to $D_w$.  This removes the intervening
excursions too.  At the starting endpoint remove the interval through
the last visit to its disk; at the terminal endpoint remove the
interval from the first visit to its disk.  These removed intervals
are disjoint along each edge.  The retained middle pieces are compact
arcs with endpoints on disk boundaries and interiors in
\[
 \Omega=\R^2\setminus\bigcup_w D_w.
\]
They are pairwise disjoint, with distinct endpoints: every possible
intersection in the original drawing was a mark and has been removed.
This construction uses the first and last visits to each \emph{closed}
disk, which exist by compactness.  It does not require an arc to cross
a round disk boundary only finitely many times.

\emph{Replacing the exterior pieces.}
The finitely many disjoint compact middle pieces admit pairwise
disjoint open neighborhoods $U_P$.  For a boundary endpoint $x$ of
one piece $P$, take a short straight spur from $x$ radially outward
from its disk.  Its interior lies in $U_P\cap\Omega$.
The spur tip can be connected to a nearby interior point of $P$ in
that same open set.  To see this directly, a sufficiently small
exterior collar at $x$ lies in $U_P$ and meets no other disk.  In polar
coordinates about the disk center, choose the collar to have radius
in $(r,r+\delta)$ and angle in a small open interval about the angle
of $x$, where $r$ is the disk radius.  It contains both the spur tip
and all sufficiently nearby interior points of $P$.  Radial paths
and a short concentric arc connect those points inside the collar.
Do this at both ends of $P$ and use the original middle subarc between
the selected interior points.  The resulting connection between the
two spur tips lies in $U_P\cap\Omega$.

Any continuous path with compact image in an open subset of the plane
can be replaced, with the same endpoints, by a polygonal path in that
open set: cover its image by small balls in the set, subdivide its
parameter interval so each subarc lies in one such ball, and join
successive subdivision points by straight segments there.  Apply this
to the connection just obtained, then attach the two spurs.  If the
result has loops, subdivide the finite union of its segments at all
segment endpoints and intersections; for collinear overlaps, also
subdivide at the overlap endpoints and identify coincident pieces.
This is a finite embedded graph.  A simple graph path between the
original boundary endpoints is a simple polygonal arc, still in
$U_P$ with its interior in $\Omega$.  Carry out this construction for
every $P$.  Their disjoint neighborhoods ensure that no resulting
exterior pieces meet.

\emph{Filling the disks.}
In a crossing disk join each of its two pairs of boundary endpoints
by a straight chord.  The four endpoints are distinct.  The chords
can therefore meet at most once, and any meeting is an interior
transverse crossing: a line cannot contain four distinct points of a
circle.  In a vertex disk choose a new vertex point in its interior
off the finitely many lines through pairs of its boundary endpoints,
and join all the endpoints to it by straight spokes.  These spokes
meet only at the vertex.  The new vertex points in different disks
are distinct.

Assemble the spokes, chords, and exterior polygonal pieces in their
original edge order.  Each resulting edge is simple: the exterior
pieces are disjoint, the disks are disjoint, and an edge uses at most
one chord or spoke in any one disk.  Its interior avoids all vertices.
Every inter-edge crossing lies in an original crossing disk, with at
most one crossing in that disk and no third edge there.  Thus the
result is admissible, polygonal, and has no more crossings than the
given drawing.  It is straight near every vertex and every crossing.

\emph{Removing adjacent-edge crossings.}
Suppose two polygonal edges $e=vu$ and $f=vw$ sharing $v$ cross at
$x$.  As the graph is simple, $u\ne w$.  Choose a small round disk
about $x$ that meets only the two straight crossing strands and
contains no other mark or bend.  Swap the portions of $e$ and $f$
from $v$ up to this disk, then connect the swapped prefixes to the
original suffixes.  Inside the disk this pairs adjacent boundary
endpoints, so two disjoint chords make the reconnections.  The result
is a walk from $v$ to $u$ and a walk from $v$ to $w$.  The crossing
at $x$ has disappeared.  Outside the disk every old open track is
assigned exactly once in total to the two new walks.

The walks may have self-intersections.  Subdivide their tracks at
all remaining intersections and bends, and repeatedly delete the
closed subwalk between two occurrences of the same subdivision
vertex.  Each deletion shortens a finite walk, so it yields a simple
polygonal path between the required endpoints.  We check that this
operation preserves admissibility.  At an old crossing of $e$ and
$f$ distinct from $x$, either the two straight-through tracks belong
to different new walks, or both belong to one walk.  In the first
case each walk visits that point only once; loop deletion either
removes its passage or retains its original straight-through
passage.  In the second case loop deletion may create a turn at that
point, but neither the other walk nor a third edge is present there.
Likewise, a crossing with any third edge occurs on a track visited
at most once by the two walks together, and any retained passage is
straight-through.  These observations cover every possible meeting:
the original drawing had no triple crossings and no common subarcs.
The tracks avoid graph vertices in their interiors, and each walk
visits its initial and terminal graph vertices only at its ends.
Hence the simplified paths are admissible edges, with no new
crossing or improper meeting.

The crossing at $x$ has been removed and none has been added.
Repeat this operation while an adjacent-edge crossing remains.  The
nonnegative integer crossing count decreases at every step, so the
process terminates.  Remaining crossings are still locally straight;
any new turn occurs away from all other edges.  The straight spokes
at vertices are merely reassigned among the two incident edges.
Small disjoint vertex, crossing, and bend disks can be chosen afresh
afterwards.  This proves all asserted properties.
\end{proof}
